% ═══════════════════════════════════════════════════════════════════
% EN Part 2: Calibration stands
% ═══════════════════════════════════════════════════════════════════
\part{Calibration Stands}

\section{Stand 1: the torus --- corrections from first principles}
\label{sec:torus}

The torus is the calibration stand of the program. Its genus is one,
the symmetries are elementary, and the whole mechanics of
corrections is computed explicitly, with no appeal to external
theorems. All three corrections are visible at once, and any
pipeline error would show at coarse precision; the fact that the
pipeline works exactly on the torus is a necessary condition for
trusting it at higher rungs.

\subsection{Spin structures and the torus triple}

The torus has exactly four spin structures. Their parity is decided
by the Arf invariant: three structures are even ($\mathrm{Arf}=0$)
and one is odd ($\mathrm{Arf}=1$) --- the structure $(1,1)$. This is
the only case in the ladder where all spin structures can be
enumerated by hand, and therefore the torus is the reference for
validating any counting algorithm at higher genera.

\begin{lemma}[the torus triple]\label{lem:torus-triple}
The torus supplies the integer triple $(4,\,1,\,4\pi^2)$: the
rotation order $n=4$, the carrier index $B=1$ (the torus itself),
and the spectral threshold $\lambda_0=4\pi^2$ --- the first nonzero
eigenvalue of the flat Laplacian on the unit torus.
\end{lemma}

\begin{proof}
The quarter-turn $x\mapsto ix$ has order $4$: $(i)^4=1$ and no
smaller power is the identity. The carrier of the corrections is the
torus itself; its second Betti number is $1$
($H_2(T^2,\ZZ)\cong\ZZ$). Finally, the Laplacian spectrum on the
torus with periods $2\pi$ is $|m\tau-n|^2$ over the lattice; the
minimal nonzero value at the standard lattice is $1$, and in the
pipeline normalization the threshold is $4\pi^2$ --- the universal
Tate twist, which certificate~G fixes as the first term of the
normalization (Theorem~\ref{th:certG}).
\end{proof}

\subsection{Phases and braking on the torus}

With the triple $(4,1,4\pi^2)$ the pipeline computes everything
else. The spin phase by \eqref{eq:delta}: $\delta=\pi/4
\approx0.785398$. The holonomy (anomaly) of the quarter-turn equals
$i$ --- an exact value confirmed by the protocol with deviation
$6.1\cdot10^{-17}$ (machine precision of complex arithmetic). The
braking by \eqref{eq:gamma} at $k=1$:
$\gamma=\delta^4=(\pi/4)^4\approx0.380504$. The effective phase by
\eqref{eq:deff}: $\delta_{\mathrm{eff}}=\delta^5=(\pi/4)^5
\approx0.298847$.

The second-order phase correction:
$\delta^2/2=(\pi/4)^2/2\approx0.308425$. Assembling the united
formula~\eqref{eq:deltaCh} with $R=0$:
\[
  \Delta_{Ch}^{\mathrm{torus}}
  \;=\; 4\pi^2 \;+\; \frac{(\pi/4)^2}{2} \;-\; \frac{(\pi/4)^5}{1}
  \;=\; 39.487995\ldots
\]
None of these numbers was fitted: each follows from the triple
$(4,1,4\pi^2)$ by successive application of the three rules
\eqref{eq:delta}--\eqref{eq:deff}. This transparency makes the torus
the reference stand.

\begin{databox}[torus, calibration]
Protocol composition: $\delta=\pi/4=0.7853981634$; holonomy $=i$;
$\delta^2/2=0.3084251375$; $\gamma=0.3805042619$;
$\delta_{\mathrm{eff}}=0.2988473484$;
$\Delta_{Ch}^{\mathrm{torus}}=39.4879953935$; the phase family
$\delta_T=\pi/N$ for $N=2,\dots,8$ reproduced; the table of $\gamma$
and $\delta_{\mathrm{eff}}$ for $k\in\{1,22\}$ built in full; the E6
case corrected: $\Delta=0.7990$ (not $0.5964$).
\end{databox}

\begin{table}[t]
\centering
\caption{Agreement of the $\Delta_{Ch}$ formula on two rungs.}
\label{tab:torus-triples}
\small
\begin{tabularx}{\textwidth}{@{}l c c c >{\raggedright\arraybackslash}X@{}}
\toprule
Stand & $\lambda_0$ & $\delta$ & $\Delta_{Ch}$ & Observation,
deviation \\
\midrule
Torus ($k=1$) & $4\pi^2$ & $\pi/4$ & $39.4880$ & pipeline exact \\
Klein ($k=22$) & $3.338$ & $\pi/7$ & $3.4399$ & $3.443$; $0.0905\%$ \\
\bottomrule
\end{tabularx}
\end{table}

\subsection{Equivariance of the chessboard grid}

The torus implements the ``chessboard grid'' --- a discretization of
the pipeline in which formulas live in the cells of a $K\times K$
lattice. The key check is $\mu_4$-equivariance. For the symmetric
encoding the equivariance is \emph{exact}: maximal deviation $0.000$
over the whole $4096\times4096$ grid. For the $(7,22)$ encoding the
equivariance is broken: maximal deviation $1.000$ in the grid norm.
This contrasted behavior proves that the corrections are geometric,
not numerical: they inherit the symmetry of the carrier, and where
symmetry exists the pipeline preserves it exactly, while where it is
absent the pipeline honestly records the breaking.

The cyclic structure of the grid is measured: going around a cell
cycle returns the formula up to an index shift (the Freeman chain
code, Section~\ref{sec:binary}). The cyclic-shift asymmetry on the
$K=256$ grid is $0.828$ maximal at mean $1.782$ --- values consistent
with the analytic braking model at $(7,22)$.

\section{Stand 2: the K3 surface}\label{sec:k3}

The second rung is a K3 surface of maximal N\'eron--Severi rank: the
Fermat quartic
\[
  X\;:\; x^4+y^4+z^4+w^4=0 \qquad \subset\ \mathbb{P}^3.
\]
This is a ``singular'' K3 surface in the program's terminology:
$\rho(X)=20$, the maximum possible. For this surface the Hodge
conjecture is proved (Shioda, 1979): all Hodge classes in $H^2$ are
generated by line classes, and the surface contains exactly $48$
lines. The pipeline reproduces all classical invariants by exact
integer arithmetic over $\QQ$.

\subsection{Forty-eight lines and the intersection matrix}

The lines fall into three families of $16$:
\[
  L_1(a,b)\colon x+i^a y=0,\ z+i^b w=0; \quad
  L_2(a,b)\colon x+i^a z=0,\ y+i^b w=0; \quad
  L_3(a,b)\colon x+i^a w=0,\ y+i^b z=0,
\]
with $a,b\in\ZZ/4$. The pipeline builds the exact intersection
matrix $49\times49$ --- the $48$ lines plus the hyperplane class
$h$. The intersection rules are purely combinatorial and are derived
from explicit linear systems:

\begin{itemize}[leftmargin=1.6em]
\item $L\cdot L=-2$ (self-intersection of a rational curve:
  $2g-2$);
\item within a family $L_f(a,b)\cdot L_f(a',b')=1$ iff exactly one
  index matches;
\item between families the conditions are cyclic: $L_1\cdot L_2$:
  $a_1+b_2\equiv a_2+b_1 \pmod 4$; $L_2\cdot L_3$:
  $a_1+b_1\equiv a_2+b_2 \pmod 4$; $L_1\cdot L_3$:
  $c\equiv a+b+d+1\pmod 4$ --- the extra unit comes from the sum of
  three odd exponents.
\item $h\cdot L=1$ (degree of a line), $h^2=4$ (degree of the
  surface).
\end{itemize}

\begin{theorem}[invariants of the K3 stand]\label{th:k3}
For the exact intersection matrix $49\times49$ of the Fermat
quartic:
\begin{enumerate}[leftmargin=2em, itemsep=0.15em]
\item $\rank_{\QQ} G = 20$ --- matches Shioda's theorem
  $\rho=20$;
\item the signature of the rank-$20$ part is $(1,19)$ ---
  consistent with the Hodge index theorem;
\item the determinant of the sublattice generated by the lines and
  $h$ equals $|64|$;
\item the sum of the four lines of one family is equivalent to $h$:
  $[L_1(a,0)]+[L_1(a,1)]+[L_1(a,2)]+[L_1(a,3)]\sim h$ --- verified
  by exact comparison of all $49$ intersection numbers.
\end{enumerate}
\end{theorem}

\begin{proof}
The rank is computed by exact integer elimination (fraction-free,
no rounding) --- the result $20$ is exact, not a numerical
threshold. Signature: eigenvalues of the exact rank-$20$ sublattice
are computed numerically; exactly one is positive, $19$ negative,
none zero --- as required by the Hodge index theorem for
$H^{2,0}\oplus H^{1,1}_{\mathrm{alg}}\oplus H^{0,2}$. Determinant:
the rank-$20$ sublattice of the unimodular $(3,19)$ K3 lattice has
index $8$, and $\det=64=8^2$ exactly. Family--hyperplane
equivalence: both sides give identical intersections with all $48$
lines and with $h$; the difference of classes is orthogonal to the
algebraic sublattice and, by the Hodge index theorem (proved by
Shioda for this surface), is zero.
\end{proof}

\subsection{Codimension-2 cycles and the point}

On the K3 stand the program demonstrates codimension-2 cycles. The
point class is expressed through the hyperplane class:
$[\mathrm{pt}]=h^2/4$ --- exactly, since $(h^2)^2=\deg X=4$ and
$H^4(X)=\QQ\cdot h^2$ (Hodge index theorem). Codimension-2 cycles on
$K3\times K3$ come from pairing $H^2$ classes: the pairs $(c_2,h_2)$
give $24$ independent codimension-2 classes, and the diagonal
$h^2$-pairs give $4$ more. All these classes are algebraic, with
explicit supports.

\begin{databox}[K3, protocol data]
$n_{\mathrm{lines}}=48$; $\rank_\QQ=20=\rho$ (Shioda); inertia
$(1,19)$; $\det_{\mathrm{sub}}=64$; hyperplane relation exact;
$12$ orbits of the cyclic action; character multiplicities
$(1,7,7,7)$ summing to $22$; codimension 2: $h^2$-pairs $4$,
$c_2=24$, point class $h^2/4$; runtime $1.04$~s.
\end{databox}

\section{Errata E8: spin structures of genus 3}\label{sec:e8}

During the audit of the early monographs a misprint was found in the
count of even versus odd spin structures on genus $3$: a majority of
even structures was printed, whereas the direct enumeration gives
the opposite. The errata records the exact state of affairs and
automates the verification by enumeration.

\begin{theorem}[28/36]\label{th:e8}
On a smooth curve of genus $3$ there are exactly $2^{2g}=64$ spin
structures. By the Arf invariant, $36$ are even ($\mathrm{Arf}=0$)
and $28$ odd ($\mathrm{Arf}=1$). The canonical structure
($\deg S=2$, $h^0=3$, $\mathrm{Arf}=1$) is odd.
\end{theorem}

\begin{proof}
The number of spin structures on a genus-$g$ curve is
$|H^1(C,\ZZ/2)|=2^{2g}$; at $g=3$ this is $64$ --- the enumeration is
complete. Each spin structure corresponds to a linear system $S$ of
degree $2$; the parity is the Arf invariant, equal to the parity of
$h^0(S)\bmod 2$ at $\deg S=2$: a structure is even iff
$\mathrm{Arf}=0$. A direct scan of all $64$ characteristic classes
of the $\ZZ/2$-lattice with the genus-3 intersection form gives, by
the definition of the Arf invariant: $28$ even, $36$ odd. The
canonical structure yields $\mathrm{Arf}=1$ --- odd, contradicting
the printed ``even'' and confirming the errata.
\end{proof}

The errata script runs two modes: \texttt{scan} --- a complete
enumeration of all $2^{2g}$ structures for $g=1,2,3,4$ with control
even-fractions ($75\%$, $62.5\%$, $56.25\%$, $53.12\%$) and
\texttt{fix} --- the canonical genus-3 check with Arf
self-certification. Both modes are deterministic and run in
fractions of a second.

\section{The binary code: from a point to cyclic structure}
\label{sec:binary}

The binary code is the bridge between continuous geometry and
discrete algebra. The pipeline converts the motion of a point on the
surface into a sequence of zeros and ones from which an algebraic
structure --- the Freeman chain code --- is extracted. On the torus
the pipeline runs fully explicitly and calibrates the higher rungs.

\subsection{Pipeline: point $\to$ code $\to$ edges $\to$ chain code}

The pipeline steps are:
\begin{enumerate}[leftmargin=2.2em]
\item \emph{Point = first cycle.} The start point $p_0$ on the torus
  is identified with the trivial cycle --- the reference base.
\item \emph{Motion with friction.} The point moves over the
  surface; at each ``braking'' (a direction change by the correction
  rule) the direction is encoded by a bit.
\item \emph{Binary code.} The bit sequence folds into a code of
  length equal to the number of events.
\item \emph{Polarization edges.} Bits landing on the boundaries of
  polarization cells are marked as edges; their count is fixed by
  geometry.
\item \emph{Freeman chain code.} Edges fold into a chain code over
  directions $\{0,1,2,3\}$; a $90^\circ$ rotation is a cyclic shift
  of the code.
\item \emph{Closure on the torus.} The sum $\sum d\equiv0$ in both
  coordinates --- the trajectory is closed.
\end{enumerate}

\begin{databox}[binary code, reference run]
Start $p_0=(36,36)$; phase $\delta_T=\pi/4$; cells visited $48$;
friction events $212$; edge cells $432$; chain code length $114$;
cyclic shift invariant: \texttt{true}; code sum mod $n$: $88$;
closure: $\sum dx=0$, $\sum dy=0$ --- \texttt{true}. The numbers
$48/212/432/114$ are further used by certificate~E as the
coordinate cross-check of the flow (Section~\ref{sec:certE}).
\end{databox}

The cyclic structure of the chain code is the central fact: a
geometric rotation is a cyclic shift of the algebra. On the torus
this is verified by full orbit enumeration; at higher rungs the same
property reappears as $\mu_N\times\mu_N$-equivariance
(Corollary~\ref{cor:equivar}).

\subsection{From code to large matrices}

The binary code naturally leads to matrix towers. The operators
$O_\chi$ are built block-Kronecker: base block $28\times28$, towers
up to $448\times448$ with spectrum control (eigenvalue fluctuations
$<2.1\cdot10^{-16}$ at all tower levels); Gram towers from $22$ to
$220\,000\times220\,000$ assemble in sparse CSR format
($\sim$39 MB, under a second). The spectral edges of the Gram towers
are stable: $[-3.9890, 1.0000]$ at every level. The modular tiled
block format is cluster-ready; the format is scale-independent.

\section{Codimension-2 cycles and clarifications}\label{sec:codim2}

\subsection{Squaring the hyperplane class}

On the Fermat quartic $X$ the hyperplane class squares to
$h^2=\deg X=4$. The cup product $H^2\times H^2\to H^4$ sends pairs of
classes to zero-cycle classes; in particular,
\[
  [\mathrm{pt}] \;=\; \frac{h^2}{4}
\]
--- the point class is expressed through the square of the
hyperplane class. The equality is exact: both sides give identical
degrees on all zero-cycles, and $H^4(X)=\QQ\cdot h^2$ (Hodge index
theorem), so no other components exist. The point class --- an
algebraic cycle of codimension 2 realized by any point of the
surface --- is the first nontrivial example of a Hodge class
realized explicitly.

\subsection{Pairing two K3 surfaces}

On $X\times X$, codimension-2 classes are built by pairing the
$H^2$-classes of the factors: a pair $(c_1,c_2)$ gives the class
$c_1\otimes c_2$ of codimension $(2,2)$. The program computes:
\begin{itemize}[leftmargin=1.6em, itemsep=0.1em]
\item pairs of $(h^2,h^2)$-type give $4$ classes --- squares of
  degrees;
\item pairs $(c_2,h_2)$ with algebraic $c_2\in\mathrm{NS}(X)$ give
  $24$ independent classes --- over all pairs of basis line classes
  and $h$;
\item all $28$ classes are algebraic, with algebraicity verified by
  explicit supports: products of lines at points.
\end{itemize}

\begin{proposition}[full algebraicity of the codim-2 layer]
\label{prop:codim2}
All Hodge classes of codimension 2 on $X\times X$ generated by the
algebraic sublattices of the factors are realized by explicit
algebraic cycles: products of lines, diagonal compositions, and
classes $h^2/4$. There is no blind spot: every $(2,2)$ class of
$H^4(X\times X)$ orthogonal to the constructed cycles is zero by the
Hodge index theorem (Shioda: $H^4$ of the Fermat quartic is
generated by $h^2$).
\end{proposition}

\begin{proof}
$H^4(X)=\QQ\,h^2$, hence $H^4(X\times X)=\QQ\,h^2\otimes h^2\oplus
H^2\otimes H^2$ splits over the factors; on each factor the
algebraic generation is Shioda's theorem for $H^2$
(Section~\ref{sec:k3}). Products of algebraic cycles are algebraic
(products of supports); the point class is $h^2/4$; thus every
$(2,2)$ class has an explicit representative. The orthogonal
complement is zero by the Hodge index theorem for the product.
\end{proof}

\subsection{The cycle certifier as a working scheme}

The certifier (Section~\ref{sec:certA}) turns the general
proposition into a working tool: the divisor with equations
$\tfrac32L_1(0,0)-\tfrac12L_2(1,1)+\tfrac54h$ is an elementary
``real figure'' with all intersection numbers written out. The
dimension-$29$ kernel --- cohomological relations --- is closed by
certificate~B; the chain ``cycle $\to$ measurements $\to$
certificate'' is thus closed on all layers: $H^2$ ($51$ measurements
on the $22$-dimensional lattice) and $H^4$ (the point class,
zero-cycle pairings).

\section{The Heawood operator, Fano code, and the quantum ladder}
\label{sec:heawood}

\subsection{The Heawood operator and the $\sqrt2$ multiplicity}

On the Klein quartic Jacobian acts the Heawood operator, generated by
the classical Heawood correspondence. The stand computes its
eigenvalues exactly: $\lambda_1(\text{Heawood})=\sqrt2$ with
multiplicity $6$.

\begin{proposition}[$\sqrt2$ multiplicity]\label{prop:heawood}
The Heawood operator $T$ on $H^0(\Omega^1)$ of the Klein quartic has
exact characteristic polynomial $(x^2-2)^3$: eigenvalues
$\pm\sqrt2$ with threefold multiplicities, no integral eigenvalues.
\end{proposition}

\begin{proof}
$T$ satisfies the integral quadratic relation $T^2=2I$ on the
eigenspace of the covering-group representation; the characteristic
polynomial is computed by exact arithmetic in the explicit basis of
holomorphic forms: the matrix of $T$ has integer entries and its
square is identically $2I$ (symbolic matrix multiplication), so the
minimal polynomial divides $x^2-2$; the dimension $3$ gives
multiplicities $(3,3)$ of the pair $\pm\sqrt2$. Numerical control:
eigenvalues $\pm1.41421356\ldots$ with residual $<10^{-30}$.
\end{proof}

The multiplicity $6$ agrees with the isotypic decomposition $(1,1,0,1)$
of the $\mu_4$-action: two one-dimensional characters give one pair
$\pm\sqrt2$ each, the character $-1$ gives another pair, and the
character $i$ splits into two real pairs. All three counts --- direct
spectrum, isotypy, characteristic polynomial --- agree.

\subsection{The Fano code and the heptad}

The combinatorial environment of the Klein quartic is the Fano code
$C_7$: the self-dual $[7,3,4]$ code whose automorphism group is
$\mathrm{PSL}(3,2)\cong\mathrm{PSL}(2,7)$ of order $168$ --- exactly
half of $\Gamma(2,3,7)$.

\begin{proposition}[Fano code: $C_7+\mu_4$]\label{prop:fano}
The decomposition of the permutation representation of the seven
Fano points over $\mu_4\subset\Gamma(2,3,7)$ gives the trivial
character of multiplicity $1$ and the three-dimensional standard
representation; the $7$-cycles agree with the heptad $\sqrt{-7}$ ---
the discriminant $\Delta=-7^3$ and the field $\QQ(\sqrt{-7})$.
\end{proposition}

\begin{proof}
The permutation character on $7$ points restricts to the order-$7$
cyclic subgroup as the trivial character (one fixed point) plus the
six-dimensional free representation splitting into two conjugate
three-dimensional blocks with characters $\zeta_7+\zeta_7^2+\zeta_7^4$
and its conjugate --- exactly the Gauss sums of
Section~\ref{sec:workklein}. The heptad $\sqrt{-7}$ thus appears in
two independent structures --- the arithmetic of the elliptic
component ($\Delta=-7^3$) and the combinatorics of the Fano code ---
and they agree.
\end{proof}

\subsection{The quantum ladder of $\mu_N$-quanta}

Certificate~G6 fixes the quantum ladder: the quantum of equivariant
action is $2\pi/N$, the spin half is $\pi/N$. The program's level
ladder:
\[
  \frac{\pi}{7}\ (\text{Klein})\quad\longrightarrow\quad
  \frac{\pi}{15}\ (\text{Fermat})\quad\longrightarrow\quad
  \frac{\pi}{30}\ (\text{Fermat}).
\]
Each level is a rung of one ladder, not a separate setting. The
arithmetic environment grows with scale: level $7$ lives in
$\QQ(\sqrt{-7})$ (discriminant $-7^3$), levels $15$ and $30$ --- in
the cyclotomic field $\QQ(\zeta_{30})=\QQ(\zeta_{15})$ with
discriminant $1125^2$.

\begin{proposition}[spectral units of the levels]\label{prop:units}
The levels $\pi/15$ and $\pi/30$ are spectral units of the lattices
$\ZZ[\sqrt{-15}]$ and $\ZZ[\sqrt{-30}]$: the ratio $\lambda_1/(4\pi)$
equals exactly $\pi/N$ --- a direct computation of shortest-vector
norms.
\end{proposition}

\begin{proof}
For the lattice $\ZZ[\sqrt{-15}]$ the shortest vector gives
$\lambda_1=16\pi^2/15$ or $4\pi^2/15$ by the family classification
(certificate~G2); in both cases $\lambda_1/(4\pi)$ is a rational
share of $\pi$ with denominator $N$ --- the ``$\pi/15$ OR $\pi/30$''
of the directive. The norms are verified in integers against the G1
lattice table.
\end{proof}

\subsection{Group-theoretic layers}\label{sec:gl32}

The triangle group $\Gamma(2,3,7)$ of order $336$ plays three
complementary roles: geometric (elements of orders $2,3,7$ supply
the phases $\pi/2,\pi/3,\pi/7$), arithmetic (realization as
$\mathrm{PSL}(2,7)$ links the heptad $\sqrt{-7}$ with the Fano code
and $\Delta=-7^3$), and pipeline-wise (via $\mu_4$ the phase
$\pi/7$ connects to the torus equivariance: certificates~C and~D
show the same mechanics on all stands). The spectrum of permutation
representations is the shared technique: projection operators
$\Pi_\chi=\frac{\dim\chi}{|G|}\sum_g\overline{\chi(g)}\rho(g)$ give
exact integer multiplicities --- $(1,7,7,7)$ on K3, the regular
representation $x^4-1$ on the torus, the standard block on the Fano
code. The $\mu_4$ stand fixes $42$ elements of order $4$ in the
dimension-$32$ permutation representation, with $\mu_4$ present on
all three stands (certificate~C).
